% a06_BALtorus_ct5: Step 3(F), first part: Ba2 p. 568, (2.46)-(2.47) (the restrictions and the Neumann cut at the
% conditional integral), with the Neumann-cut lemma. Input from sec/a06_BALtorus_c.tex.

\paragraph{Step 3(F): the remaining asserted estimates of \cite{Ba2}.} The estimates of \cite[Sect.~2]{Ba2} asserted at the conditional integral \cite[p.~568]{Ba2}, at the quadratic-form step \cite[(2.108)--(2.109)]{Ba2}, at the collar \cite[Prop.~2.1, (2.57)]{Ba2} (its field-dependent part), at the passage from \cite[(3.5)]{Ba2} to \cite[(3.21)]{Ba2}, and the region-entropy estimate \cite[(3.42)]{Ba2} are proved in this step, at block factor $5$.

\subsubsection*{The conditional integral: \cite[p.~568, (2.46)--(2.47)]{Ba2}}

\emph{Setting.} The step from scale $k-1$ to scale $k$, in scale-$(k-1)$ units, $t\defeq L^{k-1}\eps$; $\Lambda_i\defeq\Lambda^{(k-1)}_i$, $\Omega\defeq B^{k-1}(\Lambda_2)$, $\Lambda'$ the set of coarse points (the $L$-blocks of the unit lattice) whose block lies in $\Lambda$; $\varphi,A$ the new fields on coarse points, $\varphi',A'$ the old fields on unit points; $A$ as an operator argument is the background $A^{(k),\eps}$. The exponential of \cite[(2.46)]{Ba2} has ten terms: $T_1$, $T_2$ (old-field vector forms on $\Lambda^{(k-2)'}_6\cap\Lambda_5^c$ with $\Delta^{(k-1)}$ and $\Delta^{(k-1)}C^{(k-1)}_{\Lambda_5}\Delta^{(k-1)}$), $T_3=-a(L^k\eps)^{-2}\ip{(\Lambda^{(k-2)'}_6\cap\Lambda_5^c)A'}{\Delta^{(k-1)}C^{(k-1)}_{\Lambda_5}Q^*A}$, $T_4=-\frac12\ip A{\Delta^{(k),L^k\eps}_{\Lambda_5}A}$, $T_5,T_6$ (old-field forms on $\Lambda^{(k-2)'}_6\cap\Lambda_3^c$, not touched on p.~568), $T_7=-\frac12\ip{(\Lambda_3\cap\Lambda_5^c)\varphi'}{\Delta^{(k-1)}(\Omega,A)(\Lambda_3\cap\Lambda_5^c)\varphi'}$, $T_8$ the same with $\Delta C_{\Lambda_5}\Delta$, $T_9=-a(L^k\eps)^{-2}\ip{(\Lambda_3\cap\Lambda_5^c)\varphi'}{\Delta^{(k-1)}C^{(k-1)}_{\Lambda_5}Q^*(A)\varphi}$, $T_{10}=-\frac12\ip\varphi{\Delta^{(k),L^k\eps}_{\Lambda_5}(\Omega,A)\varphi}$. The coarse forms are \cite[(2.27)]{Ba2}: $\ip\psi{\Delta^{(k)}_\Lambda\psi}=aL^{d-2}\sum_{y\in\Lambda'}|\psi(y)|^2-a^2L^{-4}\ip\psi{QC^{(k-1)}_\Lambda Q^*\psi}$. \cite[p.~568]{Ba2} asserts: (2.47), the localisation of $T_{10}$ to the three forms on $\Lambda'^c_6$, $\Lambda'_6\cap\Lambda'^c_7$ and $\Lambda'_6$ with the operators on $B^{k-1}(\Lambda_2\cap\Lambda_7^c)$ (``Neumann boundary conditions'') and on $\Omega$, up to $O((L^{k-1}\eps)^\kappa)|\Lambda_6|$; the same for $T_4$ ``except that Neumann boundary conditions are not introduced'' (vector forms on the torus); the restriction of the linear terms $T_3$, $T_9$ to $\Lambda'_5\cap\Lambda'^c_6$; and the Neumann cut in $T_7$--$T_9$. Readings: $S\defeq\partial B^{k-1}(\Lambda_7)$, $\Omega_1\defeq B^{k-1}(\Lambda_2\cap\Lambda_7^c)$, $\Omega_{\mathrm{cut}}$ is $\Omega$ with the bonds crossing $S$ removed, so $\Omega_{\mathrm{cut}}=\Omega_1\sqcup B^{k-1}(\Lambda_7)$; ``Neumann conditions on $\partial B^{k-1}(\Lambda_2\cap\Lambda_7^c)$'' is $\Omega\to\Omega_{\mathrm{cut}}$, and the printed $\Delta_{\Lambda_5}(B^{k-1}(\Lambda_2\cap\Lambda_7^c),A)$ is the $\Omega_1$-block of $\Delta_{\Lambda_5}(\Omega_{\mathrm{cut}},A)$ (\cref{lem:a06f:F2}(d); read literally the printed operator is undefined, since \cite[(2.27)]{Ba2} needs $\Lambda\subset\Omega^{(k-1)}$); the third line of (2.47) reads $B^{k-1}(\Lambda_2)$; the vector forms are the instance $A=0$, $\lmu$ for $m$, componentwise. Inputs: (G) the geometry of \cite[(2.7)--(2.8)]{Ba2}: with $d_S(z)\defeq\dist(B^{k-1}(z),S)$, $z\notin\Lambda_6$ gives $d_S(z)\ge r-c_d-1$, $z\in\Lambda_3\cap\Lambda_5^c$ gives $d_S(z)\ge2(r-c_d)-1$, and $\dist(\Lambda_6,\Lambda_5^c)$, $\dist(\Lambda_6^c,\Lambda_7)\ge r-c_d$, $r=r(t)$; (FR) the field restrictions \cite[(2.16)--(2.17)]{Ba2} at scale $k-1$ \cite[p.~560]{Ba2}: $|\varphi|\le M_\varphi\defeq cp(t)\lambda(t)^{-1/4}$, $|A|\le M_A\defeq cp(t)t^{-1}$ on $\Lambda'_5$, the old fields the same on $\Lambda_3\cap\Lambda_5^c$, and the old $A'$ on $\Lambda^{(k-2)'}_6\cap\Lambda_5^c$ by the printed factor $\chi_{k-1,\Lambda^{(k-1)c}_5}$ of \cite[(2.45)--(2.46)]{Ba2}; $M\defeq\max(M_\varphi,M_A,1)\le Ct^{-2}$; (RG) $A^{(k),\eps}$ regular on $\Omega$ (\cite[p.~562, (2.98)]{Ba2}) and $e\le e_0$; (A) the propagator estimates of group (A), i.e.\ Step~3(A) at the step index $k-1$, under its hypotheses ($\Omega$ a union of big blocks, every pin set $\Lambda\subset\Omega^{(k-1)}$ a union of $L$-blocks (the proofs of Step~3(A) use of (H1) only this), $A$ regular on $\Omega$ by (RG), $e_{k-1}\le e_1$): for every pin set $\Lambda$ the operator $T_\Lambda(\Omega)\defeq-\Delta^{\eta,N}_{A,\Omega}+m^2_{k-1}+a_{k-1}P_{k-1}(\Omega\setminus B\Lambda)+a_kL^{-2}P_k(B\Lambda)$ is $\ge\gamma_L$ (\cref{lem:a06A:coer}, which is \cite[(1.8)]{Ba4} on a general region); the Combes--Thomas bound $\norm{\mathbf 1_QT^{-1}f},\norm{\mathbf 1_QD_AT^{-1}f}\le Ce^{-\theta_0\dist(Q,\supp f)}\norm f$ for unit cubes $Q$ and $T=D_A^*D_A+V+\sum_Bc_B\Pi_B\ge\gamma$ ($V\ge0$, $\Pi_B$ projections on disjoint sets of diameter $\le\ell_*$, $0\le c_B\le a$), $\theta_0=\theta_0(\gamma,\ell_*)$ (\cref{lem:a06A:CT}(i)--(ii), paired with the functions, resp.\ bond functions, supported in $Q$; $C\defeq\max\{2\gamma^{-1},(2\gamma^{-1/2}+6\sqrt d\,\theta_0\gamma^{-1})e^{\theta_0}\}$; this contains the $G_k$ clause of \cite[Cor.~2.3]{Ba4}); the decay $|\Delta^{(k-1)}(\Omega,A;z,z')|\le c_7e^{-\delta_7|z-z'|}$ (the first bound of \cref{prop:a06A:Ba1P22}, which is \cite[(2.27)]{Ba1}; its proof gives the constant $a+2a^2\gamma_\emptyset^{-1}e^{\theta_0(\gamma_\emptyset,1)}$ and the rate $\theta_0(\gamma_\emptyset,1)$, both depending on $(d,N,\norm q,a,L)$ only) and $|C_\Lambda(y,y')|\le c_2e^{-\theta_0|y-y'|}$ (\cref{lem:a06A:blockinv}(c), which is \cite[(1.16)]{Ba4}); the block inversion $X_\Lambda^{-1}=W^*G_{k-1}(\Omega,\Lambda,A)W+a_{k-1}^{-1}-(a_kL^{-2}/a_{k-1}^2)P_\Lambda$, $X_\Lambda\defeq(\Delta^{(k-1)}+aL^{-2}P)|_\Lambda$, $W\defeq Q_{k-1}^*\Lambda(1-(a_k/a_{k-1})L^{-2}P_\Lambda)$ (\cref{lem:a06A:blockinv}(a), which is \cite[(3.8)]{Ba4}); and the restriction difference $|C_\Lambda(y,y')-C^{(k-1)}(\Omega,A;y,y')|\le c_4\exp[-\frac{\theta_0}6(|y-y'|+\dist(y,\Omega^{(k-1)}\setminus\Lambda)+\dist(y',\Omega^{(k-1)}\setminus\Lambda))]$ (\cref{lem:a06A:CLdiff}, which is \cite[(1.18)]{Ba4}). On a torus these hold with the same constants (\cref{lem:a06A:torus}).

\begin{lemma}[Neumann cut]\label{lem:a06f:F1}
Let $\Omega$ be a union of big blocks with $A$ regular on $\Omega$, $Y\subset\Omega^{(k-1)}$ a union of large blocks, $S\defeq\partial B^{k-1}(Y)$, $\Omega_{\mathrm{cut}}$ as above, $\Lambda\subset\Omega^{(k-1)}$ a pin set, $T\defeq T_\Lambda(\Omega)$ and $T_{\mathrm{cut}}\defeq T_\Lambda(\Omega_{\mathrm{cut}})$ (the Neumann Laplacian of $\Omega_{\mathrm{cut}}$, same pins). Then $T_{\mathrm{cut}}\ge\gamma_L$, and for $f,f'\in\ell^2(\Omega)$ supported in sets of diameter $\le s$, both at distance $\ge R\ge4$ from $S$: $|\ip f{(T_{\mathrm{cut}}^{-1}-T^{-1})f'}|\le c_F(s)\norm f\norm{f'}\exp[-\frac{\theta_0}4(R+\dist(\supp f,\supp f'))]$, $c_F$ depending on $(d,L,a,\gamma_L,s)$ only. (For $R<4$ the bound follows from the Combes--Thomas bound applied to each resolvent.)
\end{lemma}

\begin{proof}
(1) The proof of \cref{lem:a06A:coer} drops all bonds between cells, and the cut bonds join cells since $Y$ is a union of large blocks, so $T_{\mathrm{cut}}\ge\gamma_L$; \cref{lem:a06A:CT} holds for $T_{\mathrm{cut}}$ with the bond set of $\Omega_{\mathrm{cut}}$ (its proof uses only that each site is $b_+$ of at most $d$ bonds). Equivalently: $\Omega_{\mathrm{cut}}=(\Omega\setminus B^{k-1}(Y))\sqcup B^{k-1}(Y)$ is a disjoint union of two unions of big blocks, $T_{\mathrm{cut}}=T_{\Lambda\setminus Y}(\Omega\setminus B^{k-1}(Y))\oplus T_{\Lambda\cap Y}(B^{k-1}(Y))$, and the two lemmas apply to each summand. (2) Let $\zeta(x)\defeq\min\{1,\max\{0,\frac4R\dist(x,S)-1\}\}$: $\zeta=0$ where $\dist(\cdot,S)\le R/4$, $\zeta=1$ where $\dist(\cdot,S)\ge R/2$, $|\zeta(x)-\zeta(x')|\le\frac4R|x-x'|$; so $\zeta=1$ on $\supp f\cup\supp f'$ and $\zeta$ vanishes at both ends of every cut bond (they lie within $\eta\le1\le R/4$ of $S$). (3) Put $u\defeq T^{-1}f'$; $T(\zeta u)=f'+[T,\zeta]u$. With $T-T_{\mathrm{cut}}=D_J^*D_J$, $D_J$ the covariant difference on the cut bonds, $D_J(\zeta u)=0$, so $T_{\mathrm{cut}}(\zeta u)=T(\zeta u)$ and $\zeta u=T_{\mathrm{cut}}^{-1}(f'+[T,\zeta]u)$; pairing with $f$: $\ip f{T_{\mathrm{cut}}^{-1}f'}-\ip f{T^{-1}f'}=-\ip f{T_{\mathrm{cut}}^{-1}[T,\zeta]T^{-1}f'}$. (4) With $\zeta$ acting on bond functions by $\zeta(b_-)$, $D_A(\zeta\phi)=\zeta D_A\phi+K\phi$, $(K\phi)(b)\defeq\eta^{-1}(\zeta(b_+)-\zeta(b_-))U(A_b)\phi(b_+)$, $\norm K\le4\sqrt d/R$; $[D_A^*D_A,\zeta]=D_A^*K-K^*D_A$; the potential commutes with $\zeta$; $\norm{[\Pi_B,\zeta]}\le2\operatorname{osc}_B\zeta\le8\sqrt dL/R$. Every term of $[T,\zeta]$ is supported in both variables in the layer $\mathcal L\defeq\{R/4-c\le\dist(\cdot,S)\le R/2+c\}$, $c=\sqrt dL$, and the bonds where $K\ne0$ are not cut bonds. (5) Cover $\mathcal L$ by unit cubes $Q$; the Combes--Thomas bounds for $T_{\mathrm{cut}}$ and $T$ give $\norm{\mathbf 1_QD_{A,\mathrm{cut}}T_{\mathrm{cut}}^{-1}f},\norm{\mathbf 1_QT_{\mathrm{cut}}^{-1}f}\le Ce^{-\theta_0\dist(Q,\supp f)}\norm f$ and the same for $T^{-1}f'$; writing the pairing as $\ip{D_{A,\mathrm{cut}}T_{\mathrm{cut}}^{-1}f}{KT^{-1}f'}-\ip{KT_{\mathrm{cut}}^{-1}f}{D_AT^{-1}f'}+\sum_Bc_B\ip{T_{\mathrm{cut}}^{-1}f}{[\Pi_B,\zeta]T^{-1}f'}$, each term couples neighbouring cubes only, so its modulus is $\le CR^{-1}\sum_{Q\subset\mathcal L}e^{-\theta_0(\dist(Q,\supp f)+\dist(Q,\supp f'))}\norm f\norm{f'}$. (6) For $Q\subset\mathcal L$ both distances are $\ge R/2-2c$ and their sum is $\ge\dist(\supp f,\supp f')-\sqrt d$; so $\theta_0(a+a')\ge\frac{\theta_0}4(R-4c)+\frac{\theta_0}4(\dist(\supp f,\supp f')-\sqrt d)+\frac{\theta_0}2a$, and $\sum_Qe^{-(\theta_0/2)\dist(Q,\supp f)}\le C(1+s)^d$ uniformly in the torus.
\end{proof}

\begin{lemma}[{the operators of \cite[(2.27), (2.46)]{Ba2} under the cut}]\label{lem:a06f:F2}
Take $Y=\Lambda_7$, $\Omega=B^{k-1}(\Lambda_2)$. Write $\Delta\defeq\Delta^{(k-1)}(\Omega,A)$, $\Delta^{\mathrm{cut}}\defeq\Delta^{(k-1)}(\Omega_{\mathrm{cut}},A)$; $C_\Lambda$, $C^{\mathrm{cut}}_\Lambda$ for $\Lambda\subset\Lambda_2$; $\Delta^{(k)}_\Lambda$, $\Delta^{(k),\mathrm{cut}}_\Lambda$ the coarse forms built from them, and $\Delta^{(k)}\defeq\Delta^{(k)}_{\Lambda_2}$. There are $c,\theta_*>0$, depending on $(d,N,\norm q,a,L,\gamma_L)$ only, such that for unit points $z,z'$ and coarse points $y,y'$: \textup{(a)} $|\Delta|,|\Delta^{\mathrm{cut}}|(z,z')\le ce^{-\theta_*|z-z'|}$, the same for $C_\Lambda$, $C^{\mathrm{cut}}_\Lambda$ and for the coarse forms; \textup{(b)} $|\Delta-\Delta^{\mathrm{cut}}|(z,z')\le ce^{-\theta_*(d_S(z)+d_S(z')+|z-z'|)}$, the same for $C_\Lambda-C^{\mathrm{cut}}_\Lambda$ and $\Delta^{(k)}_\Lambda-\Delta^{(k),\mathrm{cut}}_\Lambda$ (with $d_S(y)\defeq\dist(B^{k-1}B(y),S)$); \textup{(c)} for $y,y'\in\Lambda'$, $|\Delta^{(k)}_\Lambda-\Delta^{(k)}|(y,y')\le ce^{-\theta_*(|y-y'|+\dist(y,\Lambda_2\setminus\Lambda)+\dist(y',\Lambda_2\setminus\Lambda))}$; \textup{(d)} block structure: $G_{k-1}(\Omega_{\mathrm{cut}},\Lambda,A)=G_{k-1}(\Omega_1,\Lambda\cap\Lambda_7^c,A)\oplus G_{k-1}(B(\Lambda_7),\Lambda\cap\Lambda_7,A)$, and likewise $\Delta^{\mathrm{cut}}$, $C^{\mathrm{cut}}_\Lambda$, $\Delta^{(k),\mathrm{cut}}_\Lambda$ split into their $\Omega_1$- and $B(\Lambda_7)$-blocks; \textup{(e)} for products $K_1K_2K_3$ of up to three kernels of \textup{(a)} (and $Q^*$, block-local), $|K_1K_2K_3-K^{\mathrm{cut}}_1K^{\mathrm{cut}}_2K^{\mathrm{cut}}_3|(z,z')\le c'e^{-(\theta_*/3)(d_S(z)+d_S(z')+|z-z'|)}$.
\end{lemma}

\begin{proof}
(a) The first bound of \cref{prop:a06A:Ba1P22} (\cite[(2.27)]{Ba1}) for $\Delta$ and \cref{lem:a06A:blockinv}(c) (\cite[(1.16)]{Ba4}) for $C_\Lambda$: on $\Omega$ directly, and on $\Omega_{\mathrm{cut}}$ because their proofs use only \cref{lem:a06A:coer,lem:a06A:CT} and the algebra of \cref{lem:a06A:blockinv}(a), all valid on $\Omega_{\mathrm{cut}}$ by \cref{lem:a06f:F1}(1) (equivalently, by (d), whose proof is algebraic: each operator on $\Omega_{\mathrm{cut}}$ is the direct sum of its blocks on the two unions of big blocks $\Omega_1$ and $B^{k-1}(\Lambda_7)$); for the coarse kernels, $|(QKQ^*)(y,y')|\le L^d\max_{z\in B(y),z'\in B(y')}|K(z,z')|$ and the diagonal term of \cite[(2.27)]{Ba2} is local. (b) For $\Delta$: $(\Delta-\Delta^{\mathrm{cut}})(z,z')=a_{k-1}^2\ip{q(z)}{(G_{k-1}(\Omega_{\mathrm{cut}},A)-G_{k-1}(\Omega,A))q(z')}$, $q(z)\defeq Q_{k-1}^*\delta_z$ of diameter $\le\sqrt d$ and norm $1$; \cref{lem:a06f:F1} with $\Lambda=\emptyset$ and $R=\min(d_S(z),d_S(z'))$, and $\min(d_S(z),d_S(z'))+|z-z'|\ge\frac12(d_S(z)+d_S(z')+|z-z'|)-c''$. For $C_\Lambda$: by the block inversion on $\Omega$ and on $\Omega_{\mathrm{cut}}$ (algebra, needing only $T_\Lambda>0$), $W$ depends only on the pins, so $(C_\Lambda-C^{\mathrm{cut}}_\Lambda)(y,y')=\ip{W\delta_y}{(G_{k-1}(\Omega,\Lambda)-G_{k-1}(\Omega_{\mathrm{cut}},\Lambda))W\delta_{y'}}$ with $W\delta_y$ supported in the $L$-block of $y$, norm $\le1$; \cref{lem:a06f:F1} with $s=\sqrt dL$. The coarse forms follow. (c) \cref{lem:a06A:CLdiff} (\cite[(1.18)]{Ba4}; here $\Omega^{(k-1)}=\Lambda_2$, so $C^{(k-1)}(\Omega,A)=C_{\Lambda_2}$ and the distances are to $\Lambda_2\setminus\Lambda$), then the coarse step. (d) The Neumann Laplacian of $\Omega_{\mathrm{cut}}$ is the direct sum over $\Omega_1$ and $B(\Lambda_7)$; the pins are block-local and every block lies in one part, since $\Lambda_7$ is a union of large blocks; so $T_\Lambda(\Omega_{\mathrm{cut}})$ is block diagonal, $\Delta^{\mathrm{cut}}(z,z')$ vanishes unless $B(z),B(z')$ lie in the same part, $X^{\mathrm{cut}}_\Lambda$ is block diagonal and so is its inverse. (e) Telescope $K_1K_2K_3-K^{\mathrm{cut}}_1K^{\mathrm{cut}}_2K^{\mathrm{cut}}_3=(K_1-K^{\mathrm{cut}}_1)K_2K_3+K^{\mathrm{cut}}_1(K_2-K^{\mathrm{cut}}_2)K_3+K^{\mathrm{cut}}_1K^{\mathrm{cut}}_2(K_3-K^{\mathrm{cut}}_3)$; in the first term, with $X\defeq d_S(z)+d_S(w)+|z-w|$, $Y\defeq|w-w'|$, $Z\defeq|w'-z'|$ and the target $T\defeq d_S(z)+d_S(z')+|z-z'|$, the $1$-Lipschitz property of $d_S$ gives $T\le X+2Y+2Z$, so $\theta_*(X+Y+Z)\ge\frac{\theta_*}3T+\frac{\theta_*}3(Y+Z)$, and $e^{-(\theta_*/3)(Y+Z)}$ pays the sums over $w,w'$ uniformly in the torus; the other terms are the same with the roles exchanged. (A target rate $\theta_*/2$ is false: nothing would be left for the $w'$-sum.)
\end{proof}

\begin{proposition}[{\cite[p.~568]{Ba2}}]\label{prop:a06f:C10}
Assume (G), (FR), (RG), (A). For every $\kappa>0$ there is $t_0(\kappa)>0$ such that for $t=L^{k-1}\eps\le t_0$ each of the following replacements changes the exponential of \cite[(2.46)]{Ba2} by at most $Ct^\kappa|\Lambda^{(k-1)}_6|$, $C$ independent of $k$, $\eps$, the torus and the fields: \textup{(i)} $T_{10}\to$ the right side of \cite[(2.47)]{Ba2} without its $O$-term; \textup{(ii)} $T_4\to$ its analogue with $A$ in place of $\varphi$, every vector form the torus form; \textup{(iii)} $T_3\to T_3$ with $A$ replaced by $(\Lambda'_5\cap\Lambda'^c_6)A$; \textup{(iv)} $T_9\to T_9$ with $\varphi$ replaced by $(\Lambda'_5\cap\Lambda'^c_6)\varphi$ and $\Delta^{(k-1)}$, $C^{(k-1)}_{\Lambda_5}$ taken on $\Omega_{\mathrm{cut}}$, i.e.\ (by \cref{lem:a06f:F2}(d)) $\Delta^{(k-1)}(\Omega_1)$ and $C^{(k-1)}_{\Lambda_5\cap\Lambda_7^c}(\Omega_1)$; \textup{(v)} $T_7,T_8\to$ the same with every operator on $\Omega_{\mathrm{cut}}$. In particular $\kappa>d$ is allowed, as \cite[(2.48)]{Ba2} requires.
\end{proposition}

\begin{proof}
$\theta_*$ from \cref{lem:a06f:F2}; all sums are over unit or coarse points with torus-free constants. (1) Counting: for $z\notin\Lambda_6$, $d_S(z)\ge r-c$ by (G), and every unit cube meeting $S$ is adjacent to a cube of $\Lambda_7$, so for any set $X$ with $d_S\ge r-c$ on $X$, $\sum_{z\in X}e^{-\theta d_S(z)}\le e^{-\theta(r-c)/2}\sum_{Q\cap S\ne\emptyset}\sum_ze^{-\theta\dist(z,Q)/2}\le C_\theta e^{-\theta r/2}|\Lambda_7|$. (2) (i): write $\varphi=\varphi_o+\varphi_m+\varphi_i$ supported in $\Lambda'_5\setminus\Lambda'_6$, $\Lambda'_6\setminus\Lambda'_7$, $\Lambda'_7$; with $\Delta_5\defeq\Delta^{(k)}_{\Lambda_5}(\Omega)$ (symmetric), $\ip\varphi{\Delta_5\varphi}=\ip{\varphi_o}{\Delta_5\varphi_o}+2\ip{\varphi_o}{\Delta_5\varphi_m}+2\ip{\varphi_o}{\Delta_5\varphi_i}+\ip{\varphi_m+\varphi_i}{\Delta_5(\varphi_m+\varphi_i)}$; by \cref{lem:a06f:F2}(d) the first two terms of (2.47) are $\frac12\ip{\varphi_o}{\Delta^{(k),\mathrm{cut}}_{\Lambda_5}\varphi_o}+\ip{\varphi_o}{\Delta^{(k),\mathrm{cut}}_{\Lambda_5}\varphi_m}$ and the third is $\frac12\ip{\varphi_m+\varphi_i}{\Delta^{(k)}(\varphi_m+\varphi_i)}$. The difference is $E_1+E_2+E_3$: $E_1\defeq\frac12\ip{\varphi_o}{(\Delta_5-\Delta^{(k),\mathrm{cut}}_{\Lambda_5})\varphi_o}+\ip{\varphi_o}{(\Delta_5-\Delta^{(k),\mathrm{cut}}_{\Lambda_5})\varphi_m}$, by \cref{lem:a06f:F2}(b) and (1) $|E_1|\le cM^2\sum_{y\in\supp\varphi_o}e^{-\theta_*d_S(y)}\sum_{y'}e^{-\theta_*|y-y'|}\le CM^2e^{-\theta_*r/2}|\Lambda_7|$; $E_2\defeq\ip{\varphi_o}{\Delta_5\varphi_i}$, with $|y-y'|\ge r-c$ between $\supp\varphi_o$ and $\Lambda'_7$ (G), $|E_2|\le CM^2e^{-\theta_*r/2}|\Lambda_7|$; $E_3\defeq\frac12\ip{\varphi_m+\varphi_i}{(\Delta_5-\Delta^{(k)})(\varphi_m+\varphi_i)}$, by \cref{lem:a06f:F2}(c) with $\dist(y,\Lambda_2\setminus\Lambda_5)\ge r-c$ on $\Lambda_6$, $|E_3|\le CM^2e^{-\theta_*r}|\Lambda_6|$. (3) (ii): as (i) with $E_1$ absent (no cut) and the torus operators; $E_3$ by \cref{lem:a06A:CLdiff} (\cite[(1.18)]{Ba4}) on the torus, at $A=0$ with $\Omega=T_\eps$ and $\Lambda=\Lambda_5$ (\cref{lem:a06A:torus}). (4) (iii): the error is $a(L^k\eps)^{-2}\ip{A'_{\mathrm{out}}}{\Delta C_{\Lambda_5}Q^*(\Lambda'_6A)}$ with $A'_{\mathrm{out}}$ supported in $\Lambda_5^c$; by \cref{lem:a06f:F2}(a) $|(\Delta C_{\Lambda_5}Q^*)(x,y')|\le ce^{-\theta_*|x-y'|/2}$ and $|x-y'|\ge r-c$ (G), so the error is $\le Ct^{-2}M^2e^{-\theta_*r/4}|\Lambda_6|$. (5) (iv): restricting $\varphi$ is as (4); replacing $\Omega$ by $\Omega_{\mathrm{cut}}$ in $\Delta$, $C_{\Lambda_5}$ changes the kernel at $(x,y')$ by $\le ce^{-(\theta_*/3)(d_S(x)+d_S(y')+|x-y'|)}$ (\cref{lem:a06f:F2}(e)) with $d_S(x)\ge2(r-c)-1$ on $\Lambda_3\cap\Lambda_5^c$, hence by (1) $\le Ct^{-2}M^2e^{-\theta_*r/6}|\Lambda_7|$; then \cref{lem:a06f:F2}(d): $\varphi'$ is supported off $B(\Lambda_7)$ and $Q^*\varphi$ lies in $\Lambda_5\setminus\Lambda_6\subset\Lambda_7^c$, so only the $\Omega_1$-blocks act. (6) (v): as (5) with both outer points in $\Lambda_3\cap\Lambda_5^c$. (7) Every error is $\le Ct^{-4}M^2e^{-\theta_*r(t)/6}(|\Lambda_6|+|\Lambda_7|)\le C't^{-8}e^{-\theta_*r(t)/6}|\Lambda_6|$ ($M\le Ct^{-2}$, $\Lambda_7\subset\Lambda_6$), and $e^{-\theta_*r(t)/6}\le t^{\kappa+9}$ for $t\le t_0(\kappa)$ since $\varrho>1$ (the bound of Step~3(D$_2$), Geometry).
\end{proof}

The proof uses neither the Theorem of \cite{Ba4} nor the two-sided depth bounds; ``Proposition I.2.1'' of the print is not needed for these terms. The terms omitted at \cite[p.~568]{Ba2} (``described more precisely later, in (2.109)'') are the quadratic-form step treated next.
